Soit $ABC$ un triangle dont tous les angles sont aigus et tel que $AB \neq AC$. On note $\omega$ le cercle
inscrit dans $ABC$, $I$ le centre de $\omega$, et $D$, $E$ et $F$ les points de tangence respectifs de $\omega$
avec les côtés $[BC]$, $[CA]$ et $[AB]$. Soit $R$ le point de $\omega$, autre que $D$, tel que la droite $(DR)$
soit perpendiculaire à $(EF)$. Soit $P$ le point d'intersection, autre que $R$, entre la droite $(AR)$ et le
cercle $\omega$. Enfin, soit $Q$ le point d'intersection, autre que $P$, entre les cercles circonscrits à
$PCE$ et à $PBF$.

Démontrer que les droites $(DI)$ et $(PQ)$ sont sécantes en un point appartenant à la perpendiculaire à $(AI)$
passant par $A$.
